\documentclass[11pt, a4paper]{article}

\usepackage[utf8]{inputenc}
\usepackage[T1]{fontenc}
\usepackage[spanish]{babel}
\usepackage[a4paper, margin=2.5cm, top=2.5cm, bottom=2.5cm]{geometry}
\usepackage{lmodern}
\usepackage{xcolor}
\usepackage{listings}
\usepackage{amsmath}
\usepackage{booktabs}
\usepackage{graphicx}
\usepackage{caption}
\usepackage{enumitem}
\usepackage{tikz}
\usepackage{hyperref}

% Definición de Colores Institucionales y de Diseño
\definecolor{mainnavy}{RGB}{20, 45, 85}
\definecolor{secblue}{RGB}{41, 128, 185}
\definecolor{accentgold}{RGB}{212, 143, 56}
\definecolor{darkgray}{RGB}{50, 50, 50}
\definecolor{lightbg}{RGB}{248, 249, 250}
\definecolor{codebg}{RGB}{245, 247, 250}
\definecolor{codeblue}{RGB}{30, 80, 165}
\definecolor{codegreen}{RGB}{15, 125, 75}
\definecolor{codered}{RGB}{180, 40, 40}

% Estilo de Listings para C y C++
\lstdefinestyle{codigo}{
    backgroundcolor=\color{codebg},
    basicstyle=\ttfamily\footnotesize\color{darkgray},
    keywordstyle=\color{codeblue}\bfseries,
    commentstyle=\color{codegreen}\itshape,
    stringstyle=\color{codered},
    numberstyle=\tiny\color{gray},
    numbers=left,
    stepnumber=1,
    numbersep=10pt,
    showstringspaces=false,
    breaklines=true,
    frame=single,
    rulecolor=\color{secblue!30},
    tabsize=4,
    xleftmargin=15pt,
    framexleftmargin=15pt
}

\hypersetup{
    colorlinks=true,
    linkcolor=mainnavy,
    urlcolor=secblue
}

\begin{document}

% ---------------------------------------------------------
% PORTADA
% ---------------------------------------------------------
\begin{titlepage}
    \centering
    \pagecolor{white}
    
    % Barra Decorativa Superior
    \begin{tikzpicture}[remember picture, overlay]
        \fill[mainnavy] (current page.north west) rectangle ([yshift=-1.2cm]current page.north east);
        \fill[accentgold] ([yshift=-1.2cm]current page.north west) rectangle ([yshift=-1.4cm]current page.north east);
    \end{tikzpicture}
    
    \vspace*{0.8cm}
    
    % Membrete Universitario
    {\Large \textbf{\MakeUppercase{Universidad Nacional del Altiplano}}}\\[0.25cm]
    {\large \textbf{\MakeUppercase{Facultad de Ingeniería Estadística e Informática}}}\\[0.15cm]
    {\small \textbf{Escuela Profesional de Ingeniería Estadística e Informática}}\\[0.6cm]
    
    % Logos Institucionales (Verificación local limpia)
    \begin{figure}[htbp]
        \centering
        \begin{minipage}{0.4\textwidth}
            \centering
            \IfFileExists{logo-una-puno.png}{%
                \includegraphics[width=0.85\linewidth, height=3cm, keepaspectratio]{logo-una-puno.png}%
            }{%
                \framebox[0.85\linewidth][c]{\rule{0pt}{2.5cm}\small Logo UNA Puno}%
            }
        \end{minipage}
        \hspace{0.5cm}
        \begin{minipage}{0.4\textwidth}
            \centering
            \IfFileExists{logofinesi.png}{%
                \includegraphics[width=0.85\linewidth, height=3cm, keepaspectratio]{logofinesi.png}%
            }{%
                \framebox[0.85\linewidth][c]{\rule{0pt}{2.5cm}\small Logo FINESI}%
            }
        \end{minipage}
    \end{figure}
    
    \vspace{0.8cm}
    
    % Bloque del Título Principal
    \begin{tikzpicture}
        \node[fill=lightbg, draw=mainnavy, line width=1.5pt, inner sep=15pt, rounded corners=4pt, text width=0.88\textwidth, align=center] {
            {\huge \color{mainnavy} \textbf{CAPÍTULO I}}\\[0.4cm]
            {\Large \color{darkgray} \textbf{SOLUCIÓN DE EJERCICIOS DE PROGRAMACIÓN}}\\[0.2cm]
            {\large \color{secblue} \textbf{Asignatura: Estructura de Datos}}
        };
    \end{tikzpicture}
    
    \vfill
    
    % Cuadro de Información Académica y Enlace
    \begin{tikzpicture}
        \node[fill=white, draw=secblue!40, line width=1pt, inner sep=12pt, rounded corners=3pt, text width=0.88\textwidth] {
            \begin{tabular}{ll}
                \textbf{\color{mainnavy}Docente:} & Fred Torrez Cruz \\[0.15cm]
                \textbf{\color{mainnavy}Estudiante:} & Alberto Roy Chambilla Murillo \\[0.15cm]
                \textbf{\color{mainnavy}Semestre Académico:} & 2026-II \\[0.15cm]
                \textbf{\color{mainnavy}Lugar y Fecha:} & Puno, Perú --- 2026 \\[0.15cm]
                \textbf{\color{mainnavy}Repositorio GitHub:} & \href{https://github.com/albertoroy3210-creator/Ejercicios-de-estructura/blob/main/Codigo/Pdf/Latex/Ejercicios%20Cap.1.tex}{\footnotesize Ver código fuente en GitHub}
            \end{tabular}
        };
    \end{tikzpicture}
    
    \vspace{0.8cm}
    
    % Barra Decorativa Inferior
    \begin{tikzpicture}[remember picture, overlay]
        \fill[accentgold] ([yshift=1.2cm]current page.south west) rectangle ([yshift=1.4cm]current page.south east);
        \fill[mainnavy] (current page.south west) rectangle ([yshift=1.2cm]current page.south east);
    \end{tikzpicture}

\end{titlepage}

% Reset de la página
\nopagecolor
\clearpage
\pagenumbering{arabic}

% ---------------------------------------------------------
% PRESENTACIÓN E INTRODUCCIÓN
% ---------------------------------------------------------
\section*{\color{mainnavy} Presentación}
El presente documento contiene la solución detallada y estructurada de los ejercicios prácticos del \textbf{Capítulo I} para el curso de \textbf{Estructura de Datos}. Cada ejercicio incluye la justificación metodológica, la implementación del código fuente en lenguaje C/C++ y la referencia a su correspondiente ejecución gráfica.

Puedes acceder al código fuente completo y actualizado en el siguiente repositorio de GitHub:
\begin{center}
    \url{https://github.com/albertoroy3210-creator/Ejercicios-de-estructura/blob/main/Codigo/Pdf/Latex/Ejercicios%20Cap.1.tex}
\end{center}

\vspace{0.3cm}
\hrulefill
\vspace{0.3cm}

% ---------------------------------------------------------
% EJERCICIOS
% ---------------------------------------------------------

\section*{\color{mainnavy}Ejercicio 1. Costo total del procesamiento}
\textbf{Procedimiento:} Se leen las horas trabajadas y el costo por hora. El costo total se obtiene multiplicando ambos valores.
\begin{lstlisting}[style=codigo,language=C]
#include <stdio.h>

int main(void) {
    float horas, tarifa, costo;

    printf("Ingrese las horas trabajadas: ");
    scanf("%f", &horas);

    printf("Ingrese el costo por hora: ");
    scanf("%f", &tarifa);

    costo = horas * tarifa;

    printf("\n--- RESUMEN ---\n");
    printf("Horas trabajadas: %.2f\n", horas);
    printf("Tarifa por hora: S/ %.2f\n", tarifa);
    printf("Costo total: S/ %.2f\n", costo);

    return 0;
}
\end{lstlisting}
\begin{figure}[htbp]
    \centering
    \IfFileExists{1.png}{%
        \includegraphics[width=0.6\linewidth, height=6cm, keepaspectratio]{1.png}%
    }{%
        \framebox[0.6\linewidth][c]{\rule{0pt}{4cm}\small Captura de pantalla Ejercicio 1 (1.png)}%
    }
    \caption{Resultado del Ejercicio 1}
    \label{fig:ej1}
\end{figure}

\vspace{0.5cm}

\section*{\color{mainnavy}Ejercicio 2. Registro de una observación}
\textbf{Procedimiento:} Se utilizan tipos de datos apropiados: entero para el identificador y la edad, float para el promedio, char para la categoría y un entero para representar el estado de validez.
\begin{lstlisting}[style=codigo,language=C]
#include <stdio.h>

int main(void) {
    int identificador, edad, valido;
    float promedio;
    char categoria;

    printf("Identificador: ");
    scanf("%d", &identificador);

    printf("Edad: ");
    scanf("%d", &edad);

    printf("Valor promedio: ");
    scanf("%f", &promedio);

    printf("Categoria (una letra): ");
    scanf(" %c", &categoria);

    printf("Registro valido? (1=Si, 0=No): ");
    scanf("%d", &valido);

    printf("\n--- DATOS DE LA OBSERVACION ---\n");
    printf("Identificador: %d\n", identificador);
    printf("Edad: %d\n", edad);
    printf("Promedio: %.2f\n", promedio);
    printf("Categoria: %c\n", categoria);
    printf("Estado: %s\n", valido == 1 ? "Valido" : "No valido");

    return 0;
}
\end{lstlisting}
\begin{figure}[htbp]
    \centering
    \IfFileExists{2.png}{%
        \includegraphics[width=0.6\linewidth, height=6cm, keepaspectratio]{2.png}%
    }{%
        \framebox[0.6\linewidth][c]{\rule{0pt}{4cm}\small Captura de pantalla Ejercicio 2 (2.png)}%
    }
    \caption{Resultado del Ejercicio 2}
    \label{fig:ej2}
\end{figure}

\vspace{0.5cm}

\section*{\color{mainnavy}Ejercicio 3. Conversión de decimal a entero}
\textbf{Procedimiento:} La conversión explícita mediante \texttt{(int)} elimina la parte decimal, sin redondear. La pérdida se calcula restando el entero convertido al valor original.
\begin{lstlisting}[style=codigo,language=C]
#include <stdio.h>

int main(void) {
    float medicion, perdida;
    int valorEntero;

    printf("Ingrese la medicion decimal: ");
    scanf("%f", &medicion);

    valorEntero = (int)medicion;
    perdida = medicion - valorEntero;

    printf("Valor original: %.2f\n", medicion);
    printf("Valor entero: %d\n", valorEntero);
    printf("Parte decimal perdida: %.2f\n", perdida);

    return 0;
}
\end{lstlisting}
\begin{figure}[htbp]
    \centering
    \IfFileExists{3.png}{%
        \includegraphics[width=0.6\linewidth, height=6cm, keepaspectratio]{3.png}%
    }{%
        \framebox[0.6\linewidth][c]{\rule{0pt}{4cm}\small Captura de pantalla Ejercicio 3 (3.png)}%
    }
    \caption{Resultado del Ejercicio 3}
    \label{fig:ej3}
\end{figure}

\vspace{0.5cm}

\section*{\color{mainnavy}Ejercicio 4. Operaciones aritméticas}
\textbf{Procedimiento:} Se leen dos enteros y se calculan suma, diferencia, producto, división entera y residuo. Se comprueba que el segundo número no sea cero antes de dividir.
\begin{lstlisting}[style=codigo,language=C]
#include <stdio.h>

int main(void) {
    int a, b;

    printf("Ingrese el primer entero: ");
    scanf("%d", &a);

    printf("Ingrese el segundo entero: ");
    scanf("%d", &b);

    printf("\nSuma: %d\n", a + b);
    printf("Diferencia: %d\n", a - b);
    printf("Producto: %d\n", a * b);

    if (b != 0) {
        printf("Division entera: %d\n", a / b);
        printf("Residuo: %d\n", a % b);
    } else {
        printf("No se puede dividir entre cero.\n");
    }

    return 0;
}
\end{lstlisting}
\begin{figure}[htbp]
    \centering
    \IfFileExists{4.png}{%
        \includegraphics[width=0.6\linewidth, height=6cm, keepaspectratio]{4.png}%
    }{%
        \framebox[0.6\linewidth][c]{\rule{0pt}{4cm}\small Captura de pantalla Ejercicio 4 (4.png)}%
    }
    \caption{Resultado del Ejercicio 4}
    \label{fig:ej4}
\end{figure}

\vspace{0.5cm}

\section*{\color{mainnavy}Ejercicio 5. Distribución de registros entre nodos}
\textbf{Procedimiento:} La división entera determina los registros asignados a cada nodo y el operador módulo determina los registros sobrantes.
\begin{lstlisting}[style=codigo,language=C]
#include <stdio.h>

int main(void) {
    int registros, nodos;

    printf("Numero total de registros: ");
    scanf("%d", &registros);

    printf("Numero de nodos: ");
    scanf("%d", &nodos);

    if (registros >= 0 && nodos > 0) {
        printf("Registros por nodo: %d\n", registros / nodos);
        printf("Registros sin distribuir: %d\n", registros % nodos);
    } else {
        printf("Ingrese registros no negativos y nodos mayores que cero.\n");
    }

    return 0;
}
\end{lstlisting}
\begin{figure}[htbp]
    \centering
    \IfFileExists{5.png}{%
        \includegraphics[width=0.6\linewidth, height=6cm, keepaspectratio]{5.png}%
    }{%
        \framebox[0.6\linewidth][c]{\rule{0pt}{4cm}\small Captura de pantalla Ejercicio 5 (5.png)}%
    }
    \caption{Resultado del Ejercicio 5}
    \label{fig:ej5}
\end{figure}

\vspace{0.5cm}

\section*{\color{mainnavy}Ejercicio 6. Operadores de incremento y asignación}
\textbf{Procedimiento:} La variable inicia en 4, luego aumenta en 3 y finalmente se multiplica por 2.
\begin{lstlisting}[style=codigo,language=C]
#include <stdio.h>

int main(void) {
    int registros = 4;

    printf("Valor inicial: %d\n", registros);

    registros += 3;
    printf("Despues de incrementar en 3: %d\n", registros);

    registros *= 2;
    printf("Despues de duplicar: %d\n", registros);

    return 0;
}
\end{lstlisting}
\begin{figure}[htbp]
    \centering
    \IfFileExists{6.png}{%
        \includegraphics[width=0.6\linewidth, height=6cm, keepaspectratio]{6.png}%
    }{%
        \framebox[0.6\linewidth][c]{\rule{0pt}{4cm}\small Captura de pantalla Ejercicio 6 (6.png)}%
    }
    \caption{Resultado del Ejercicio 6}
    \label{fig:ej6}
\end{figure}

\vspace{0.5cm}

\section*{\color{mainnavy}Ejercicio 7. Validación de un conjunto de datos}
\textbf{Procedimiento:} El conjunto se acepta únicamente si tiene al menos 100 observaciones válidas y un porcentaje de datos completos igual o superior al 70\%. Ambas condiciones deben cumplirse.
\begin{lstlisting}[style=codigo,language=C]
#include <stdio.h>

int main(void) {
    int observaciones;
    float porcentaje;

    printf("Numero de observaciones validas: ");
    scanf("%d", &observaciones);

    printf("Porcentaje de datos completos: ");
    scanf("%f", &porcentaje);

    if (observaciones >= 100 && porcentaje >= 70.0 && porcentaje <= 100.0) {
        printf("El conjunto de datos puede ser analizado.\n");
    } else {
        printf("El conjunto de datos no cumple las condiciones.\n");
    }

    return 0;
}
\end{lstlisting}
\begin{figure}[htbp]
    \centering
    \IfFileExists{7.png}{%
        \includegraphics[width=0.6\linewidth, height=6cm, keepaspectratio]{7.png}%
    }{%
        \framebox[0.6\linewidth][c]{\rule{0pt}{4cm}\small Captura de pantalla Ejercicio 7 (7.png)}%
    }
    \caption{Resultado del Ejercicio 7}
    \label{fig:ej7}
\end{figure}

\vspace{0.5cm}

\section*{\color{mainnavy}Ejercicio 8. Validación de una observación}
\textbf{Procedimiento:} Se utiliza el operador lógico AND (\texttt{\&\&}). La observación se incorpora si la edad es como mínimo 18 y el registro está validado.
\begin{lstlisting}[style=codigo,language=C]
#include <stdio.h>

int main(void) {
    int edad, validado;

    printf("Ingrese la edad: ");
    scanf("%d", &edad);

    printf("Registro validado? (1=Si, 0=No): ");
    scanf("%d", &validado);

    if (edad >= 18 && validado == 1) {
        printf("La observacion puede incorporarse.\n");
    } else {
        printf("La observacion no puede incorporarse.\n");
    }

    return 0;
}
\end{lstlisting}
\begin{figure}[htbp]
    \centering
    \IfFileExists{8.png}{%
        \includegraphics[width=0.6\linewidth, height=6cm, keepaspectratio]{8.png}%
    }{%
        \framebox[0.6\linewidth][c]{\rule{0pt}{4cm}\small Captura de pantalla Ejercicio 8 (8.png)}%
    }
    \caption{Resultado del Ejercicio 8}
    \label{fig:ej8}
\end{figure}

\vspace{0.5cm}

\section*{\color{mainnavy}Ejercicio 9. Clasificación de temperatura}
\textbf{Procedimiento:} Se utilizan condiciones encadenadas para clasificar la temperatura en tres categorías.
\begin{lstlisting}[style=codigo,language=C]
#include <stdio.h>

int main(void) {
    float temperatura;

    printf("Ingrese la temperatura en grados Celsius: ");
    scanf("%f", &temperatura);

    if (temperatura < 0) {
        printf("Clasificacion: Congelacion\n");
    } else if (temperatura <= 20) {
        printf("Clasificacion: Frio\n");
    } else {
        printf("Clasificacion: Templado\n");
    }

    return 0;
}
\end{lstlisting}
\begin{figure}[htbp]
    \centering
    \IfFileExists{9.png}{%
        \includegraphics[width=0.6\linewidth, height=6cm, keepaspectratio]{9.png}%
    }{%
        \framebox[0.6\linewidth][c]{\rule{0pt}{4cm}\small Captura de pantalla Ejercicio 9 (9.png)}%
    }
    \caption{Resultado del Ejercicio 9}
    \label{fig:ej9}
\end{figure}

\vspace{0.5cm}

\section*{\color{mainnavy}Ejercicio 10. Media y observaciones superiores a ella (C++)}
\textbf{Procedimiento:} Se almacenan diez mediciones en un arreglo y se calcula su suma. La media se obtiene dividiendo la suma entre diez. Después se cuentan los valores mayores que la media.
\begin{lstlisting}[style=codigo,language=C++]
#include <iostream>
using namespace std;

int main() {
    double datos[10], suma = 0, media;
    int contador = 0;

    for (int i = 0; i < 10; i++) {
        cout << "Ingrese la medicion " << i + 1 << ": ";
        cin >> datos[i];
        suma += datos[i];
    }

    media = suma / 10;

    for (int i = 0; i < 10; i++) {
        if (datos[i] > media) {
            contador++;
        }
    }

    cout << "Media aritmetica: " << media << endl;
    cout << "Observaciones por encima de la media: " << contador << endl;

    return 0;
}
\end{lstlisting}
\begin{figure}[htbp]
    \centering
    \IfFileExists{10.png}{%
        \includegraphics[width=0.6\linewidth, height=6cm, keepaspectratio]{10.png}%
    }{%
        \framebox[0.6\linewidth][c]{\rule{0pt}{4cm}\small Captura de pantalla Ejercicio 10 (10.png)}%
    }
    \caption{Resultado del Ejercicio 10}
    \label{fig:ej10}
\end{figure}

\vspace{0.5cm}

\section*{\color{mainnavy}Ejercicio 11. Validación de contraseña (C++)}
\textbf{Procedimiento:} El ciclo \texttt{while} solicita la clave hasta que sea correcta. Se incrementa el contador cada vez que se realiza un intento. Para este ejemplo, la clave correcta se establece en 1234.
\begin{lstlisting}[style=codigo,language=C++]
#include <iostream>
using namespace std;

int main() {
    const int CLAVE_CORRECTA = 1234;
    int clave, intentos = 0;

    cout << "Ingrese la clave numerica: ";
    cin >> clave;
    intentos++;

    while (clave != CLAVE_CORRECTA) {
        cout << "Clave incorrecta. Intente nuevamente.\n";
        cout << "Ingrese la clave numerica: ";
        cin >> clave;
        intentos++;
    }

    cout << "Acceso permitido.\n";
    cout << "Numero de intentos: " << intentos << endl;

    return 0;
}
\end{lstlisting}
\begin{figure}[htbp]
    \centering
    \IfFileExists{11.png}{%
        \includegraphics[width=0.6\linewidth, height=6cm, keepaspectratio]{11.png}%
    }{%
        \framebox[0.6\linewidth][c]{\rule{0pt}{4cm}\small Captura de pantalla Ejercicio 11 (11.png)}%
    }
    \caption{Resultado del Ejercicio 11}
    \label{fig:ej11}
\end{figure}

\vspace{0.5cm}

\section*{\color{mainnavy}Ejercicio 12. Procesamiento de mediciones con break y continue (C++)}
\textbf{Procedimiento:} Se leen como máximo diez valores. Los negativos se omiten con \texttt{continue}; el valor 999 detiene el ciclo con \texttt{break}. Solo se cuentan las mediciones no negativas distintas de 999.
\begin{lstlisting}[style=codigo,language=C++]
#include <iostream>
using namespace std;

int main() {
    double valor;
    int validos = 0;

    for (int i = 0; i < 10; i++) {
        cout << "Ingrese la medicion " << i + 1 << " (999 para terminar): ";
        cin >> valor;

        if (valor == 999) break;

        if (valor < 0) {
            cout << "Valor negativo: se omite.\n";
            continue;
        }

        validos++;
    }

    cout << "Valores validos procesados: " << validos << endl;

    return 0;
}
\end{lstlisting}
\begin{figure}[htbp]
    \centering
    \IfFileExists{12.png}{%
        \includegraphics[width=0.6\linewidth, height=6cm, keepaspectratio]{12.png}%
    }{%
        \framebox[0.6\linewidth][c]{\rule{0pt}{4cm}\small Captura de pantalla Ejercicio 12 (12.png)}%
    }
    \caption{Resultado del Ejercicio 12}
    \label{fig:ej12}
\end{figure}

\vspace{0.5cm}

\section*{\color{mainnavy}Ejercicio 13. Análisis de edades mediante un arreglo}
\textbf{Procedimiento:} Se almacenan diez edades. Durante el ingreso se determina la suma, la edad mínima y la máxima. Después se calcula la media y se cuentan las edades superiores a ella.
\begin{lstlisting}[style=codigo,language=C]
#include <stdio.h>

int main(void) {
    int edades[10], suma = 0;
    int minima, maxima, superiores = 0;
    float media;

    for (int i = 0; i < 10; i++) {
        printf("Edad del participante %d: ", i + 1);
        scanf("%d", &edades[i]);

        if (i == 0) {
            minima = maxima = edades[i];
        } else {
            if (edades[i] < minima) minima = edades[i];
            if (edades[i] > maxima) maxima = edades[i];
        }

        suma += edades[i];
    }

    media = (float)suma / 10;

    for (int i = 0; i < 10; i++) {
        if (edades[i] > media) superiores++;
    }

    printf("\nEdad minima: %d\n", minima);
    printf("Edad maxima: %d\n", maxima);
    printf("Media de edades: %.2f\n", media);
    printf("Participantes por encima de la media: %d\n", superiores);

    return 0;
}
\end{lstlisting}
\begin{figure}[htbp]
    \centering
    \IfFileExists{13.png}{%
        \includegraphics[width=0.6\linewidth, height=6cm, keepaspectratio]{13.png}%
    }{%
        \framebox[0.6\linewidth][c]{\rule{0pt}{4cm}\small Captura de pantalla Ejercicio 13 (13.png)}%
    }
    \caption{Resultado del Ejercicio 13}
    \label{fig:ej13}
\end{figure}

\vspace{0.5cm}

\section*{\color{mainnavy}Ejercicio 14. Comparación de complejidad $O(n)$ y $O(n^2)$}
\textbf{Procedimiento:} El primer procedimiento ejecuta una operación por cada elemento ($n$ operaciones). El segundo ejecuta dos ciclos anidados de $n$ iteraciones ($n^2$ operaciones).
\begin{lstlisting}[style=codigo,language=C]
#include <stdio.h>

int main(void) {
    int n;
    long long operacionesLineales = 0;
    long long operacionesCuadraticas = 0;

    printf("Ingrese el numero de elementos: ");
    scanf("%d", &n);

    if (n <= 0) {
        printf("n debe ser un entero positivo.\n");
        return 1;
    }

    /* Procedimiento 1: un recorrido */
    for (int i = 0; i < n; i++) {
        operacionesLineales++;
    }

    /* Procedimiento 2: dos ciclos anidados */
    for (int i = 0; i < n; i++) {
        for (int j = 0; j < n; j++) {
            operacionesCuadraticas++;
        }
    }

    printf("Operaciones del primer procedimiento: %lld\n", operacionesLineales);
    printf("Complejidad aproximada: O(n)\n");

    printf("Operaciones del segundo procedimiento: %lld\n", operacionesCuadraticas);
    printf("Complejidad aproximada: O(n^2)\n");

    return 0;
}
\end{lstlisting}
\begin{figure}[htbp]
    \centering
    \IfFileExists{14.png}{%
        \includegraphics[width=0.6\linewidth, height=6cm, keepaspectratio]{14.png}%
    }{%
        \framebox[0.6\linewidth][c]{\rule{0pt}{4cm}\small Captura de pantalla Ejercicio 14 (14.png)}%
    }
    \caption{Resultado del Ejercicio 14}
    \label{fig:ej14}
\end{figure}

\vspace{0.5cm}

\section*{\color{mainnavy}Ejercicio 15. Funciones para promedio y conteo}
\textbf{Procedimiento:} La función \texttt{calcular\_promedio} recibe tres mediciones y retorna su media. La función \texttt{contar\_superiores} recibe las mediciones y el promedio, y retorna cuántas superan dicho promedio.
\begin{lstlisting}[style=codigo,language=C]
#include <stdio.h>

float calcular_promedio(float a, float b, float c) {
    return (a + b + c) / 3.0f;
}

int contar_superiores(float a, float b, float c, float promedio) {
    int contador = 0;
    if (a > promedio) contador++;
    if (b > promedio) contador++;
    if (c > promedio) contador++;
    return contador;
}

int main(void) {
    float m1, m2, m3, promedio;
    int cantidad;

    printf("Ingrese la primera medicion: ");
    scanf("%f", &m1);

    printf("Ingrese la segunda medicion: ");
    scanf("%f", &m2);

    printf("Ingrese la tercera medicion: ");
    scanf("%f", &m3);

    promedio = calcular_promedio(m1, m2, m3);
    cantidad = contar_superiores(m1, m2, m3, promedio);

    printf("Promedio: %.2f\n", promedio);
    printf("Mediciones por encima del promedio: %d\n", cantidad);

    return 0;
}
\end{lstlisting}
\begin{figure}[htbp]
    \centering
    \IfFileExists{15.png}{%
        \includegraphics[width=0.6\linewidth, height=6cm, keepaspectratio]{15.png}%
    }{%
        \framebox[0.6\linewidth][c]{\rule{0pt}{4cm}\small Captura de pantalla Ejercicio 15 (15.png)}%
    }
    \caption{Resultado del Ejercicio 15}
    \label{fig:ej15}
\end{figure}

% ---------------------------------------------------------
% CONCLUSIÓN
% ---------------------------------------------------------
\vspace{0.8cm}
\section*{\color{mainnavy} Conclusión}
Los ejercicios compilados en este informe abarcan los conceptos esenciales del paradigma de programación estructurada: control de flujo mediante estructuras condicionales y bucles, tipos abstractos de datos primitivos, procesamiento secuencial y matricial con arreglos, funciones como unidades modulares y principios básicos sobre el comportamiento temporal de los algoritmos ($O(n)$ y $O(n^2)$).

\end{document}